\documentclass[UTF8,a4paper,11pt]{ctexart}
\usepackage[margin=25mm,headheight=14pt]{geometry}
\usepackage{xcolor}
\usepackage{booktabs}
\usepackage{longtable}
\usepackage{array}
\usepackage{enumitem}
\usepackage{microtype}
\usepackage[numbers,sort&compress]{natbib}
\usepackage[hidelinks]{hyperref}

\definecolor{ink}{HTML}{141414}
\definecolor{vermilion}{HTML}{D45438}
\definecolor{muted}{HTML}{6D6961}
\definecolor{paper}{HTML}{F4F1EA}
\pagecolor{paper}
\color{ink}

\setCJKmainfont{STSong}
\setCJKsansfont{Heiti SC}
\setmainfont{Georgia}
\setsansfont{Helvetica Neue}
\setmonofont{Menlo}

\ctexset{
  section = {format = \Large\bfseries\color{vermilion}, name = {}, afterskip = 8pt},
  subsection = {format = \large\bfseries\color{ink}, afterskip = 4pt},
  subsubsection = {format = \normalsize\bfseries\color{muted}}
}

\newenvironment{voicebox}{%
  \begin{quote}\small\setlength{\parskip}{5pt}\color{ink}%
}{%
  \end{quote}%
}

\newcommand{\metric}[2]{\texttt{#1}: #2}
\newcommand{\pastiche}{\textit{教学仿写（teaching pastiche）}}


\title{一题十一声：把文风从感觉变成可校验的分布}
\author{style-mimic teaching edition}
\date{\today}

\begin{document}
\maketitle
\begin{abstract}
本文提出一个面向文风分析与教学仿写的轻量框架。它把文风拆成词汇、句法、修辞、衔接语篇、立场介入和叙事声音六层，要求画像同时记录作者特征与平均作者基线的差分。我们进一步加入 p05/p95、IQR、MAD、偏度和 CV 等尾部统计，并以模式簇、内容重叠和样本门槛约束 review 结论。为展示“内容不变、声纹可变”，本文将同一语义骨架写成十一种中文与英文教学仿写版本。结果表明，文风可以被量化为可复核的选择分布，但任何启发式分数都不能替代人工阅读或作者身份判断。
\end{abstract}

\section{研究问题}

“像某位作者”通常被当成一种直觉判断，然而直觉很难告诉写作者下一步该改哪一句。更可操作的目标是：固定论点与体裁零件，只改变词汇、句法节奏、修辞、衔接、立场和叙事视角，再观察这些选择是否落入目标作者的分布区间。这个目标既比“完全复制声音”谦逊，也比“套用几个口头禅”严格。

本文的核心主张是：文风不是覆盖在内容上的滤镜，而是多层语言选择共同形成的分布。一个可用的工具必须让特征可以从原文复算，让短样本和零基线显式降级，让相似度不被单一指标垄断，并把活人作者模仿清楚标成教学练习。

\section{框架与方法}

\subsection{六层画像}

词汇层记录词汇丰富度、实词与功能词、术语和高频内容词；句法层记录句长分布、每句小句数、句类与节奏；修辞层记录明喻、反复、排比和设问；衔接语篇层记录连接标记、句间词汇重叠与段落推进；立场介入层记录 hedge、booster、作者自称、读者代词和指令；叙事声音层则记录 POV、话语呈现和聚焦。这种从词汇、句法到叙事声音的分层，承接了经典文体学的可观察特征传统 \citep{leech1981style,biber1988variation,hyland2005stance}。前五层回答“怎样说”，第六层回答“谁在说、谁在看”。

\subsection{基线与尾部}

若只记录“某作者常用因此”，画像很容易把文体常识误报成个人风格。因此我们先构造同题目的平均作者基线，再用目标作者减去基线，保留差分特征。校验时也不能只看平均句长：两个文本可能拥有相同均值，却一个长短交错、另一个句长被熨平。为此，分析报告同时输出中位数、p05/p95、四分位距 IQR、中位绝对偏差 MAD、偏度与变异系数 CV；MATTR 还可降低文本长度差异对词汇丰富度比较的影响 \citep{covington2010mattr}；少于 20 句时，尾部只作线索。

\subsection{Review 而非判决}

\texttt{aiflavor} 把 AI 味拆成可追溯的模式簇：空洞大词、升华套话、模糊归因、套路开头结尾、排比模式与节奏过匀，并显示原文证据。\texttt{sim} 将六层声纹距离和内容 Jaccard 分开，避免两个词汇完全不重叠的文本因为 TTR 相同而获得虚假的高相似度。这个设计也回应了近期研究对语言变异度、模型平均化与多信号评测的提醒 \citep{zanotto2025linguistic,reinhart2025llm,zanotto2024variability}。少于 5 句或 80 个单位时，系统只报告线索，不下“同一声纹”结论。

\section{共享语义骨架与多声部案例}

为了控制变量，我们把论文固定成摘要、方法、案例、限制与结论五个单元。十一种版本都保留以下事实：style-mimic 使用六层画像；画像需要平均作者差分；校验要看分布尾部；review 需要多信号；仿写必须标明伦理边界。变化只发生在表达层：有的版本用游记般的旷达，有的版本用白话实验，有的版本用演讲的重复，有的版本用科学类比。

下面的章节由 `data/styles.json` 自动生成。每个框中的文本都属于 \pastiche，不是作者原作。

\input{generated-voices.tex}

\section{量化附录}

% Generated from data/styles.json; values are review signals.
\begin{longtable}{@{}lllll@{}}
\toprule Voice & Language & CV & AI flavor & Similarity \\
\midrule
苏轼 & 中文 & 0.388 & 22 & 60 \\
桐城派 & 中文 & 0.617 & 0 & 72 \\
王勃 & 中文 & 0.348 & 25 & 64 \\
鲁迅 & 中文 & 0.539 & 38 & 55 \\
胡适 & 中文 & 0.515 & 6 & 55 \\
钱钟书 & 中文 & 0.38 & 27 & 35 \\
Barbara Oakley & English & 0.509 & 14 & 48 \\
Donald Trump & English & 0.717 & 0 & 54 \\
Barack Obama & English & 0.454 & 6 & 52 \\
Albert Einstein & English & 0.447 & 22 & 39 \\
Richard Feynman & English & 0.339 & 15 & 47 \\
\bottomrule
\end{longtable}


表中的数值来自仓库内的 \texttt{style\_analyze.py} 与 \texttt{style\_review.py}。它们用于比较和改稿定位，不用于识别作者身份，也不构成“文本由 AI 生成”的结论。由于中文与英文分词方式不同、体裁也不同，跨语言分数只能作低置信度参考。

\section{结论与边界}

同一命题可以拥有许多声部，而声部差异不必靠神秘直觉来说明。把内容固定，把选择拆层，把尾部保留，把证据留在原文旁边，写作者就获得了一条可重复的练习回路：建立画像，写出变体，比较偏差，再人工朗读。这个回路不会制造“真正的作者”，但能让学习者更清楚地知道自己改变了什么。

最后必须保留三条边界。第一，短文本的统计不稳定，阈值不能脱离样本量和体裁解释。第二，AI 味是编辑线索，不是法证检测。第三，活人作者的可识别风格只能用于明确标注的教学练习，不得冒充原作者、伪造出处或制造背书。

\bibliographystyle{plainnat}
\bibliography{references}

\end{document}
